% =====================================================================
% Sekcja 7 - Ograniczenia
% =====================================================================

\section{Ograniczenia}
\label{sec:limitations}

Wyniki tej pracy trzeba czytać z kilkoma zastrzeżeniami.

Najpoważniejsze dotyczy dekodowania. Wszystkie eksperymenty
korzystają z dekodowania zachłannego, przez co model na każdym kroku
wybiera swój najbardziej prawdopodobny token, a wartości
$1/\mathrm{PPL}$ skupiają się blisko 1,0. Sygnał pewności jest więc
ściśnięty i część informacji, którą niosłoby próbkowanie z pełnego
rozkładu, zostaje utracona. Powtórzenie eksperymentów z włączonym
samplowaniem mogłoby dać bogatszy obraz - i prawdopodobnie inne,
być może niższe ECE.

Ograniczona jest też skala. Badano dwa małe modele z jednej rodziny
(Qwen2.5 0,5B i 1,5B) na jednym zbiorze - GSM8K. To zadania
arytmetyczne z dokładnym dopasowaniem odpowiedzi liczbowej, więc nie
wiadomo, na ile obserwacje przenoszą się na większe modele ani na
zadania o innym charakterze (otwarte odpowiedzi, kod, dłuższe
trajektorie z wieloma narzędziami). Agent korzystał z jednego
narzędzia - kalkulatora - co jest najprostszym możliwym ustawieniem
ReAct.

Niewielkie są wreszcie próby. Korelacje liczone są na 60 (a w pilocie
150) obserwacjach, a w analizie segmentów \texttt{Action} efektywne
$n$ spada do 48. Przy takich liczebnościach przedziały ufności są
szerokie - bootstrapowy 95\% CI dla najmocniejszej korelacji sięga
od $-0{,}60$ do $-0{,}24$ - więc raportowane współczynniki należy
traktować jako rzędy wielkości i kierunki zależności, a nie jako
precyzyjne oszacowania. Z tego samego powodu w pracy nie testowano
rekalibracji pewności; wskazano ją jedynie jako naturalny dalszy krok.
